\section{Exact Verification Suite and Reproducibility Manifest}
\label{sec:appendix}

All mathematical identities, spectral decompositions, exhaustive $S_9$ and order-$4$ censuses, continuous kernel trajectories, $n = 3\dots 50$ multi-order scaling censuses, and $n = 46, 48, 100$ head-to-head solver benchmarks reported in this paper are certified by the standalone Python verification suite \texttt{verify\_paper\_d.py}, which outputs the machine-readable JSON dataset \texttt{results\_paper\_d.json}. First, \texttt{verify\_section2\_spectral()} constructs the $(2n+2)\times n^2$ incidence matrix $A_n$ and operator $L_n = A_n^\top A_n$ for all $3 \le n \le 8$, verifies the irrational eigenvalues $\lambda_\pm(n)$ on $\R^{n^2}$ (\Cref{eq:unprojected_spectrum}), projects $L_n$ onto an orthonormal Householder basis of $\mathbf{1}^\perp$ to verify the exact integer spectrum $\spec(\Ltilde_n|_{\mathbf{1}^\perp})$ (\Cref{eq:rationalized_spectrum}), certifies $\tr(\Ltilde_n) = 2(n^2-1)$ and $\mathbb{E}[\Hline(S)] = 2\|\mathbf{z}\|^2 = \frac{2}{3}n^2(n^4-1)$ (\Cref{prop:spectral_rationalization}), and checks all residuals and moments of the $4\times 4$ low-mode counterexample $S_{\mathrm{low}}$ (\Cref{rem:low_mode_insufficiency}). Second, \texttt{verify\_section3\_multiscale()} evaluates $\widetilde{H}_p(k)$ ($p\in\{1,2,4\}$, $1\le k\le n$), interior and toroidal $2\times 2$ quartet energies $W_{2\times 2}^{\mathrm{int}}$ and $W_{2\times 2}^{\mathrm{tor}}$, complement-pair break counts $B_{\mathrm{comp}}$, the Concentric Shell Quadrupole Identity $M_{xyz}(S) = -8\,Q_{\mathrm{corner}}(S) - 3\,Q_{\mathrm{edge}}(S)$, and mixed octupole moments $M_{xyz}$ across all $880$ Fr\'enicle order-$4$ magic squares (partitioned by Dudeney Groups~I--XII) and across constructive families at $n \in \{3,4,5,6,8\}$, certifying \Cref{lem:quartet_equivalence} and \Cref{prop:octupole_census}. Third, \texttt{verify\_section4\_swap\_landscape()} evaluates $\Hline(S)$ and all $36$ transposition deltas $\Delta\Hline(\tau_{ab})$ via \Cref{lem:swap_law} across all $9! = 362{,}880$ permutations of $S_9$, verifies that the $8$ Lo Shu states and the $8$ antipodal traps $D_4\cdot S_{\mathrm{trap}}^{(3)}$ are the unique $2$-cycle local minima in $S_9$, confirms that all $168$ directed $3$-cycles strictly increase $\Hline$ on $S_{\mathrm{trap}}^{(3)}$, proves that the row-preserving double swap $(6\leftrightarrow 7)(3\leftrightarrow 4)$ is the unique escaping $(2,2)$ move among all $378$ disjoint double swaps (crossing the intermediate single-swap saddle at $\Hline^\ddagger = 12$ and restoring the steepest-descent basin of $S_9$ from $270{,}442$ to $362{,}880 / 362{,}880$, \Cref{thm:s9_traps}), and runs the $2{,}000$-trial steepest-descent, $3$-cycle, and line-preserving $(2,2)$-cycle escape census on $S_{16}$. Fourth, \texttt{verify\_section5\_bubble\_trinity()} integrates the continuous $2\mathrm{D}$ spatial kernel equations (\Cref{eq:langevin_sde}) across Regimes~I--IV (evaluating both greedy cell-snapping and exact $O(N^3)$ Hungarian minimum-weight bipartite matching \cite{kuhn1955hungarian,munkres1957algorithms}), benchmarks $(2+3)$-exchange simulated annealing against multi-exchange local search, performs the controlled $6$-tier operator ablation study (\Cref{tab:physical_analogs_ablation}), and certifies $100.0\%$ multi-order synthesis across $n \in \{3\dots 10\}$ (\Cref{tab:multi_order_settler}). Fifth, \texttt{section6\_large\_scale\_competition} in \texttt{results\_paper\_d.json} archives the complete order-by-order scaling census across all $48$ consecutive orders $n = 3\dots 50$ (\Cref{fig:scan_3_to_50} and \Cref{tab:scan_milestone_orders}), verifies the Bilinear Plaquette Identity $\Delta\Hline(p_1, p_2) = \Delta\Hline(p_1) + \Delta\Hline(p_2) + 2\,\boldsymbol{\delta R}(p_1)^\top \boldsymbol{\delta R}(p_2)$ and Active-Set Steepest-Descent Theorem (\Cref{thm:bilinear_plaquette}), and records the head-to-head solver competition across Simulated Annealing, Defective-Line CBLS, Baseline All-Pairs ME-ILS, and Bilinear-Identity + Active-Set ME-ILS at orders $n = 46$, $n = 48$, and $n = 100$ (\Cref{fig:competition_100} and \Cref{tab:head_to_head_competition}).
